% Generated draft: requires mathematical and source review.
\documentclass[11pt]{article}
\usepackage[margin=25mm]{geometry}
\usepackage{fontspec}
\setmainfont{Latin Modern Roman}
\usepackage{amsmath,amssymb,mathtools,unicode-math}
\setmathfont{Latin Modern Math}
\usepackage{enumitem,needspace,xurl,hyperref}
\setlength{\parindent}{0pt}
\setlength{\parskip}{5pt}
\setlength{\emergencystretch}{3em}
\newcommand{\N}{\mathbb{N}}
\newcommand{\Z}{\mathbb{Z}}
\newcommand{\Q}{\mathbb{Q}}
\newcommand{\R}{\mathbb{R}}
\newcommand{\C}{\mathbb{C}}
\newcommand{\F}{\mathbb{F}}
\newcommand{\A}{\mathbb{A}}
\newcommand{\PP}{\mathbb{P}}
\newcommand{\E}{\mathbb{E}}
\newcommand{\eps}{\varepsilon}
\newcommand{\dd}{\,\mathrm{d}}
\begin{document}
\section*{UnsolvedMath ID: 30005703}
\textbf{Source problem code:} catalogue entry
\subsection{Definitions and mathematical statement}
For an abstract group \(K\), a unital ring \(R\), and an integer \(n\geq 1\), say that \(K\) is of type \(\mathrm{FP}_n(R)\) if the trivial left \(RK\)-module \(R\) admits a projective resolution
\[
\cdots\longrightarrow P_{n+1}\longrightarrow P_n\longrightarrow\cdots\longrightarrow P_0\longrightarrow R\longrightarrow 0
\]
in which \(P_i\) is finitely generated for every \(0\leq i\leq n\). A group \(H\) is of type \(F\) if it admits a finite classifying CW complex \(K(H,1)\). Here type \(F\) denotes a geometric finiteness property and is distinct from the coefficient field denoted below by \(F\).

Put
\[
\mathcal K=\{\Q\}\cup\{\F_p:p\text{ is prime}\}.
\]
The supplied question asks for classes of groups in which finiteness over the indicated fields forces integral finiteness. Its coefficient quantifier is unresolved: the phrase ``for \(F\in\mathcal K\)'' could require finiteness for a single chosen field, or simultaneously for every field in \(\mathcal K\). These yield the following distinct formulations for a class \(\mathcal C\) and an integer \(n\geq1\):
\[
\forall F\in\mathcal K\ \forall G\in\mathcal C,\qquad
G\text{ is of type }\mathrm{FP}_n(F)
\ \Longrightarrow\
G\text{ is of type }\mathrm{FP}_n(\Z),
\tag{A}
\]
and
\[
\forall G\in\mathcal C,\qquad
\bigl(\forall F\in\mathcal K,\ G\text{ is of type }\mathrm{FP}_n(F)\bigr)
\ \Longrightarrow\
G\text{ is of type }\mathrm{FP}_n(\Z).
\tag{B}
\]
\textbf{Classification question.} For which classes \(\mathcal C\) does the intended implication hold? The supplied record does not resolve whether formulation (A), formulation (B), or a statement restricted to a specified field is intended, nor whether the requested classes may depend on \(n\) or on the field.

\textbf{Exact-sequence question.} Does the intended implication hold under the additional hypothesis that there is an exact sequence, written in the source as
\[
G\xrightarrow{\iota}H\xrightarrow{\pi}\Z,
\qquad H\text{ of type }F?
\]
Exactness at \(H\) means \(\operatorname{im}\iota=\ker\pi\). The displayed sequence supplies no endpoint terms, so injectivity of \(\iota\) and surjectivity of \(\pi\) are not explicit. If the intended hypothesis is instead the short exact sequence
\[
1\longrightarrow G\xrightarrow{\iota}H\xrightarrow{\pi}\Z\longrightarrow 1,
\]
then \(G\) is identified with the kernel of an epimorphism from a group of type \(F\) onto \(\Z\); this stronger interpretation requires source confirmation.

\textbf{Source qualifications.} The original excerpt supplies the range \(n\geq1\), asserts that the integral conclusion fails in general, and mentions right-angled Artin groups as examples where it holds. These are supplied contextual claims whose precise coefficient hypotheses are not established by the excerpt. The definitions above describe ordinary abstract-group finiteness properties. Although the workshop concerns totally disconnected locally compact groups, the supplied target does not specify a topology or a topological module category, so a topological analogue cannot be substituted without further hypotheses.
\subsection{Short English statement}
Determine when homological finiteness over rational or prime fields forces the corresponding integral finiteness property, especially for groups occurring in the displayed exact sequence with a group of type \(F\). The field quantifier, the sequence endpoints, and the intended abstract or topological setting require clarification. This is a statement-only draft requiring human mathematical review.
\subsection{Sources}
\begin{itemize}
\item Supplied local UnsolvedMath registry record, permanent ID \(30005703\), title: Integral Finiteness from Rational or Modular Finiteness; exact source problem code: \texttt{catalogue entry}. No proposer is identified for this target.
\item Supplied source citation: Mini-Workshop: Homological aspects of TDLC-groups (2024), Oberwolfach Reports, DOI 10.4171/owr/2023/55, pp.~3103--3138. Supplied source URL: \url{https://doi.org/10.4171/owr/2023/55}. The record separately gives proposed year 2023; the citation's year and the DOI's year component are retained without reconciliation.
\item Supplied literature source, labeled Oberwolfach report: \url{https://ems.press/content/serial-article-files/48166}.
\item Literature assessment dated exactly \texttt{2026-08-22}: the supplied curation labels the problem open, reports that the source poses both questions as unresolved, and reports that no later result establishing the stated implication in general was located. This is supplied dated curation, not an independently verified literature assessment.
\item The registry labels the statement \texttt{corrected\_verified} and says it was checked against the official report. That verification claim is supplied metadata; no independent source visit or verification is claimed here.
\item The original excerpt also contains a separately headed discussion attributed to Rudradip Biswas and an incomplete reference to [1, Theorem 19.1]. No bibliographic identification of [1] is supplied, and this adjacent attribution does not establish authorship of the present target.
\end{itemize}
\end{document}
